\documentclass[11pt,a4paper]{article}
\usepackage[margin=2.2cm]{geometry}
\usepackage{fontspec}
\usepackage{xeCJK}
\usepackage{booktabs}
\usepackage{tabularx}
\usepackage{longtable}
\usepackage{array}
\usepackage{enumitem}
\usepackage{hyperref}
\usepackage{xcolor}

% Font configurations
\setmainfont{DejaVu Sans}
\setCJKmainfont{Noto Sans CJK SC}

\hypersetup{
    colorlinks=true,
    linkcolor=blue,
    filecolor=magenta,
    urlcolor=blue,
    pdftitle={Ban Dich Tai Lieu Nghien Cuu Hoc Thuat - Classroom Mind Wandering},
}

\definecolor{navyblue}{rgb}{0.0, 0.2, 0.5}
\definecolor{darkblue}{rgb}{0.1, 0.1, 0.4}

\title{\textbf{BẢN DỊCH TÀI LIỆU PHỤC VỤ NGHIÊN CỨU HỌC THUẬT}\\
\large\textit{Nghiên cứu về Hiện tượng Xao nhãng / Mất tập trung trong Giờ học của Sinh viên (Mind Wandering Study)}}
\author{\textbf{Hồ sơ Dịch thuật \& Định dạng LaTeX}}
\date{\today}

\begin{document}

\maketitle

\begin{abstract}
Tài liệu này chứa toàn bộ nội dung dịch thuật đối chiếu từ tiếng Trung / tiếng Anh sang tiếng Việt phục vụ cho mục đích nghiên cứu học thuật về hiện tượng xao nhãng trong giờ học (\textit{Classroom Mind Wandering}). Dữ liệu bao gồm 25 phát biểu Q-Sort, quy trình hướng dẫn từng bước trên công cụ Easy HtmlQ, các thông số cấu hình và câu hỏi điều tra nhân khẩu học. Văn bản được biên soạn và đổ lại bằng XeLaTeX để đảm bảo chuẩn hóa định dạng, không bị lỗi font hay phá vỡ cấu trúc hiển thị.
\end{abstract}

\tableofcontents
\newpage

\section{Tổng quan và Mục đích Sử dụng}
Khảo sát này sử dụng phương pháp Q-Methodology (thông qua công cụ Easy HtmlQ) để đánh giá thái độ, hành vi và các yếu tố ảnh hưởng đến sự mất tập trung của sinh viên/lưu học sinh trong môi trường lớp học.
\begin{itemize}
    \item \textbf{Phạm vi dịch thuật}: Dịch thuật phục vụ nghiên cứu và giảng dạy học thuật (Academic Use).
    \item \textbf{Cấu trúc tài liệu}: Bao gồm danh mục 25 phát biểu Q-Sort, các thông điệp giao diện người dùng (UI) và bảng hỏi thông tin nền.
    \item \textbf{Định dạng mã hóa}: Toàn bộ file cấu hình XML đã được cập nhật song ngữ tương ứng trong thư mục \texttt{easy-htmlq-2.0.3/settings/}.
\end{itemize}

\section{Danh mục 25 Phát biểu Q-Sort (Q-Sort Statements)}

Dưới đây là bảng đối chiếu chi tiết 25 phát biểu được sử dụng trong quy trình phân loại Q-Sort:

\begin{longtable}{>{\bfseries}c p{6.5cm} p{8cm}}
\toprule
STT & Nguyên bản tiếng Trung (Chinese) & Bản dịch tiếng Việt (Vietnamese Translation) \\
\midrule
\endhead

1 & 教室中的座位安排会影响我集中注意力的能力。 & Sự sắp xếp chỗ ngồi trong lớp học ảnh hưởng đến khả năng tập trung của tôi. \\
\midrule
2 & 如果课堂管理不严格，我往往会做与学习无关的事情。 & Nếu quản lý lớp học không nghiêm ngặt, tôi thường làm những việc không liên quan đến học tập. \\
\midrule
3 & 当课程内容缺乏吸引力时，我就会分心做别的事情。 & Khi nội dung bài học thiếu sức hấp dẫn, tôi sẽ xao nhãng làm việc khác. \\
\midrule
4 & 涉及文化敏感话题的课程让我感到不适，因此我会选择回避。 & Những bài học đề cập đến các chủ đề nhạy cảm về văn hóa khiến tôi cảm thấy khó chịu, vì vậy tôi chọn cách né tránh. \\
\midrule
5 & 课程太难我听不懂时，我会失去继续听讲的动力。 & Khi bài học quá khó và không thể hiểu, tôi mất động lực tiếp tục lắng nghe. \\
\midrule
6 & 课程太简单时，我认为不听也能懂，因此容易分神。 & Khi bài học quá dễ, tôi nghĩ không nghe cũng hiểu nên dễ bị phân tâm. \\
\midrule
7 & 身体不适时，我很难维持课堂注意力。 & Khi cơ thể không khỏe, tôi rất khó duy trì sự chú ý trong giờ học. \\
\midrule
8 & 只要手机或电子设备在身边，我的注意力就容易被吸走。 & Chỉ cần có điện thoại hoặc thiết bị điện tử bên cạnh, sự chú ý của tôi rất dễ bị cuốn đi. \\
\midrule
9 & 缺乏学习动机让我在课堂上难以投入。 & Thiếu động lực học tập khiến tôi khó chú tâm vào bài giảng trên lớp. \\
\midrule
10 & 同学的言行时常打断我的注意力，使我分心。 & Lời nói và hành động của bạn học thường xuyên ngắt quãng sự tập trung của tôi, làm tôi phân tâm. \\
\midrule
11 & 有时我看着老师好像在听，其实心早就飘走了。 & Đôi khi tôi nhìn thầy cô có vẻ như đang nghe, nhưng thực ra tâm trí đã bay đi xa. \\
\midrule
12 & 一旦开始走神，我的坐姿也会变得慵懒随意。 & Một khi bắt đầu xao nhãng, tư thế ngồi của tôi cũng trở nên uể uể, tùy tiện. \\
\midrule
13 & 上课时，我常陷入发呆状态，无法专注于课堂内容。 & Trong giờ học, tôi thường rơi vào trạng thái thẫn thờ/ngơ ngác, không thể tập trung vào nội dung bài giảng. \\
\midrule
14 & 有时我会用笔或小玩具在课堂上打发时间。 & Đôi khi tôi dùng bút hoặc đồ chơi nhỏ để giết thời gian trong lớp. \\
\midrule
15 & 我会在课堂上随手乱画，以此来打发无聊时间。 & Tôi hay vẽ nguệch ngoạc trong giờ học để giết thời gian nhàm chán. \\
\midrule
16 & 有时候我甚至意识不到自己已经在发呆了。 & Có những lúc tôi thậm chí không nhận ra bản thân mình đang thẫn thờ/phân tâm. \\
\midrule
17 & 即使意识到自己在走神，我也很难重新集中注意力。 & Dù nhận ra mình đang xao nhãng, tôi cũng rất khó tập trung lại sự chú ý. \\
\midrule
18 & 我的注意力经常在课堂上游离，脑子里想些无关的事情。 & Sự chú ý của tôi thường xuyên lang thang trong giờ học, đầu óc suy nghĩ về những việc không liên quan. \\
\midrule
19 & 一旦发现自己分心，我会主动采取措施进行调整。 & Ngay khi phát hiện mình phân tâm, tôi sẽ chủ động thực hiện các biện pháp để điều chỉnh lại. \\
\midrule
20 & 有些时候我就是不想听课，只想放空自己。 & Có những lúc tôi đơn giản là không muốn nghe giảng, chỉ muốn để tâm trí trống rỗng. \\
\midrule
21 & 走神之后，我常常会感到懊悔和内疚。 & Sau khi xao nhãng, tôi thường cảm thấy hối hận và cắn rứt lương tâm. \\
\midrule
22 & 我认为走神是难以避免的现象，不必过度苛责自己。 & Tôi cho rằng xao nhãng là hiện tượng khó tránh khỏi, không cần phải quá khắt khe trách phạt bản thân. \\
\midrule
23 & 当走神影响了我的学习效果，我会感到非常困惑和懊恼。 & Khi việc xao nhãng ảnh hưởng đến kết quả học tập, tôi cảm thấy rất bối rối và bực bội. \\
\midrule
24 & 如果我觉得这门课对我未来没用，就容易不专心。 & Nếu tôi cảm thấy môn học này không có ích cho tương lai của mình, tôi sẽ dễ mất tập trung. \\
\midrule
25 & 有时候，我觉得自己是在逃避课堂带来的压力。 & Đôi khi, tôi cảm thấy mình đang trốn tránh áp lực do lớp học mang lại. \\
\bottomrule
\end{longtable}

\newpage
\section{Hướng dẫn và Chuỗi Giao diện Nối tiếp (Language \& UI Items)}

\begin{description}[style=unboxed,leftmargin=0cm]
    \item[Chào mừng (Welcome Screen)] \hfill \\
    \textit{Tiếng Trung}: 这是一个关于在课堂上走神的研究项目。重要信息：在本次调查中，您需要尽可能多地使用屏幕空间！如果有必要，请将您的浏览器窗口最大化，重新加载此网页，并点击“继续”按钮开始调查。\\
    \textit{Tiếng Việt}: Đây là dự án nghiên cứu về hiện tượng xao nhãng / mất tập trung (mind wandering) của sinh viên trong giờ học. Thông tin quan trọng: Trong khảo sát này, bạn cần sử dụng tối đa không gian màn hình! Nếu cần thiết, vui lòng phóng to tối đa cửa sổ trình duyệt, tải lại trang web này và nhấp vào nút "Tiếp tục" để bắt đầu khảo sát.

    \item[Bước 1 / 5 (Step 1 of 5: Rough Sorting)] \hfill \\
    \textit{Tiếng Trung}: 请仔细阅读以下陈述，并将其分成三类：一类是您不太认同的陈述，一类是您较为赞同的陈述，还有一类是其余的陈述。您可以将这些卡片拖入这三类中的任意一类，或者在键盘上按 1、2、3 键。之后可以再做修改。如果您想再次阅读此说明，请点击左下角的帮助按钮。\\
    \textit{Tiếng Việt}: Vui lòng đọc kỹ các phát biểu sau và phân loại chúng thành ba nhóm: nhóm không đồng ý, nhóm tương đối đồng ý, và nhóm các phát biểu còn lại (trung lập). Bạn có thể kéo các thẻ này vào bất kỳ nhóm nào trong ba nhóm, hoặc nhấn phím 1, 2, 3 trên bàn phím. Sau đó bạn vẫn có thể chỉnh sửa lại. Nếu muốn đọc lại hướng dẫn này, vui lòng nhấp vào nút trợ giúp ở góc dưới bên trái.

    \item[Bước 2 / 5 (Step 2 of 5: Grid Placement)] \hfill \\
    \textit{Tiếng Trung}: 从“同意”这一组卡片中取出卡片，并将它们排列在得分表的右侧。接下来，从“不同意”这一组卡片中取出卡片，并将其排列在得分表的左侧。按照此步骤对“同意”和“不同意”这两组卡片中的所有卡片进行处理。最后，取出“中立”卡片，并将它们放置在得分表剩余的空白格中。\\
    \textit{Tiếng Việt}: Rút các thẻ từ nhóm "Đồng ý" và xếp chúng vào phía bên phải của bảng điểm. Tiếp theo, rút các thẻ từ nhóm "Không đồng ý" và xếp chúng vào phía bên trái của bảng điểm. Thực hiện theo bước này đối với tất cả các thẻ trong hai nhóm "Đồng ý" và "Không đồng ý". Cuối cùng, rút các thẻ "Trung lập" và đặt chúng vào các ô trống còn lại trên bảng điểm.

    \item[Bước 3 / 5 (Step 3 of 5: Review Sorting)] \hfill \\
    \textit{Tiếng Trung}: 现在您已经将所有卡片都摆放在计分表上了。请再仔细检查一下您的分配情况，如果需要的话可以重新调整卡片的位置。\\
    \textit{Tiếng Việt}: Bây giờ bạn đã sắp xếp tất cả các thẻ lên bảng điểm. Vui lòng kiểm tra lại cẩn thận sự phân bố của bạn, và nếu cần thiết có thể điều chỉnh lại vị trí các thẻ.

    \item[Bước 4 / 5 (Step 4 of 5: Comments)] \hfill \\
    \textit{Tiếng Trung}: 请解释一下您为何对以下陈述中的哪一条最为赞同或最为反对。请在“+4”或“-4”处标注您的意见。\\
    \textit{Tiếng Việt}: Vui lòng giải thích lý do tại sao bạn đồng ý nhất hoặc phản đối nhất đối với các phát biểu bên dưới. Vui lòng ghi lại ý kiến của bạn ở vị trí "+4" hoặc "-4".

    \item[Bước 5 / 5 (Step 5 of 5: Background Questions)] \hfill \\
    \textit{Tiếng Trung}: 最后，请回答以下关于您个人背景的问题。\\
    \textit{Tiếng Việt}: Cuối cùng, vui lòng trả lời các câu hỏi sau về thông tin cá nhân của bạn.
\end{description}

\section{Bảng hỏi Thông tin Nhân khẩu học (Demographic Survey Form)}

\begin{table}[h!]
\centering
\small
\begin{tabularx}{\textwidth}{>{\bfseries}l X X}
\toprule
Trường thông tin & Câu hỏi / Ghi chú nguyên bản & Bản dịch tiếng Việt \\
\midrule
Tuổi (Age) & 请输入您的年龄 (YYYY, eg. 2003) & Vui lòng nhập năm sinh của bạn (YYYY, ví dụ: 2003) \\
\addlinespace
Thời gian học tiếng Trung & 请输入您开始学中文的时间 (YYYY, eg. 2020) & Vui lòng nhập thời điểm bạn bắt đầu học tiếng Trung (YYYY, ví dụ: 2020) \\
\addlinespace
Trình độ HSK & 请输入您的HSK等级 (N, eg. 6) & Vui lòng nhập cấp độ HSK của bạn (N, ví dụ: 6) \\
\bottomrule
\end{tabularx}
\caption{Các câu hỏi thu thập thông tin nền của sinh viên}
\end{table}

\section{Kết luận}
Bản dịch đã được tích hợp hoàn toàn vào hệ thống Easy HtmlQ. Việc định dạng lại bằng XeLaTeX giúp giữ nguyên tính toàn vẹn của văn bản học thuật, hỗ trợ xuất bản và lưu trữ hồ sơ nghiên cứu sạch đẹp, không bị lỗi font tiếng Việt hay ký tự tiếng Trung.

\end{document}
